\item Dos puntos $A$ y $B$ sobre la superficie de la Tierra están a la misma longitud y separados $60.0^\circ$ en latitud, como se muestra en la figura P16.2. Suponga que un terremoto en el punto $A$ crea una onda P que llega al punto $B$ viajando en línea recta a través del cuerpo de la Tierra con una rapidez constante de 7.80 km/s. El terremoto también genera una onda Rayleigh que viaja a 4.50 km/s. Además de las ondas P y S, las ondas Rayleigh son un tercer tipo de onda sísmica que viaja sobre la superficie de la Tierra en lugar de hacerlo a través de su volumen.
\begin{enumerate}
    \item ¿Cuál de estas dos ondas sísmicas llega primero a $B$?
    \item ¿Cuál es la diferencia de tiempo entre las llegadas de estas dos ondas a $B$?
\end{enumerate}